% ============================================================
% Phụ lục B — Bảng các luật suy diễn áp dụng vào ontology Chèo
%
% Yêu cầu preamble:
%   \usepackage{longtable}
% ============================================================

\section*{B. Bảng các luật suy diễn áp dụng vào ontology Chèo}
\label{appendix:inference_rules}
\addcontentsline{toc}{section}{B. Các luật suy diễn áp dụng vào ontology Chèo}

\begin{longtable}{|p{0.9cm}|p{3.0cm}|p{4.0cm}|p{6.4cm}|}
  \hline
  \textbf{ID} & \textbf{Tên luật} & \textbf{Mô tả} & \textbf{Luật} \\
  \hline
  \endfirsthead
  \hline
  \textbf{ID} & \textbf{Tên luật} & \textbf{Mô tả} & \textbf{Luật} \\
  \hline
  \endhead
  \hline
  \endfoot
  \hline
  \endlastfoot

  \texttt{A.1}
  & Cross-Gender Casting Detection
  & Phát hiện vai diễn có giới tính diễn viên khác giới tính nhân vật.
  & \texttt{Actor(?a), performedBy(?ra,?a),
    forCharacter(?ra,?c), actorGender(?a,?ag),
    charGender(?c,?cg), notEqual(?ag,?cg)
    => CrossGenderCasting(?ra)} \\
  \hline

  \texttt{A.2}
  & Incomplete Character Detection
  & Phát hiện Character không thuộc phân lớp Đào/Hề/Kép/Mụ/Lão.
  & \texttt{Character(?c), NOT EXISTS \{?c rdf:type ?sub.
    FILTER(?sub IN (DaoCharacter, HeCharacter, KepCharacter,
    MuCharacter, LaoCharacter))\}
    => IncompleteCharacter(?c)} \\
  \hline

  \texttt{A.3}
  & RoleType-Gender Mismatch
  & Phát hiện sai lệch quy ước giới: Đào/Mụ phải là Nữ, Kép/Lão phải là Nam.
  & \texttt{DaoCharacter(?c), charGender(?c,"Nam") => AtypicalRole(?c)}\par
    \texttt{KepCharacter(?c), charGender(?c,"Nữ") => AtypicalRole(?c)}\par
    \texttt{MuCharacter(?c), charGender(?c,"Nam") => AtypicalRole(?c)}\par
    \texttt{LaoCharacter(?c), charGender(?c,"Nữ") => AtypicalRole(?c)} \\
  \hline

  \texttt{A.4}
  & Orphaned Appearance Detection
  & Phát hiện Appearance không gắn với RoleAssignment nào.
  & \texttt{Appearance(?app), NOT EXISTS \{?ra hasAppearance ?app\}
    => OrphanedAppearance(?app)} \\
  \hline

  \texttt{B.1}
  & Symmetric Closure
  & Đóng kín đối xứng cho ba quan hệ ngang hàng.
  & \texttt{collaboratesWith(?a,?b) => collaboratesWith(?b,?a)}\par
    \texttt{hasRelation(?a,?b) => hasRelation(?b,?a)}\par
    \texttt{isOpponentOf(?a,?b) => isOpponentOf(?b,?a)} \\
  \hline

  \texttt{B.2}
  & Direct Performs Relation
  & Rút gọn chuỗi Actor--RoleAssignment--Character thành quan hệ trực tiếp.
  & \texttt{Actor(?a), performedBy(?ra,?a),
    forCharacter(?ra,?c) => performs(?a,?c)} \\
  \hline

  \texttt{B.3}
  & CollaboratesWith from Co-Version
  & Hai diễn viên cùng tham gia một phiên bản trình diễn thì có quan hệ cộng tác.
  & \texttt{performedBy(?ra1,?a1), inVersion(?ra1,?v),
    performedBy(?ra2,?a2), inVersion(?ra2,?v),
    notEqual(?a1,?a2) => collaboratesWith(?a1,?a2)} \\
  \hline

  \texttt{B.4}
  & Versatile Actor Classification
  & Diễn viên đóng từ hai loại vai khác nhau trở lên là đa năng.
  & \texttt{Actor(?a), COUNT(DISTINCT roleType played by ?a) >= 2
    => versatility(?a, "Đa năng")}\par
    \textit{(Cần aggregation; SPARQL CONSTRUCT với GROUP BY ?a
    HAVING COUNT(DISTINCT ?role) >= 2.)} \\
  \hline

  \texttt{B.5}
  & Recurring Character Classification
  & Nhân vật xuất hiện trong $\geq 2$ trích đoạn cùng vở chèo là xuyên suốt.
  & \texttt{Character(?c), Play(?p),
    COUNT(DISTINCT scene of ?c in ?p) >= 2
    => recurring(?c, true)}\par
    \textit{(Cần aggregation; SPARQL với GROUP BY ?c
    HAVING COUNT(DISTINCT ?s) >= 2.)} \\
  \hline

  \texttt{B.6}
  & Emotional Profile
  & Tính cảm xúc chủ đạo và dải cảm xúc cho mỗi nhân vật.
  & \texttt{Character(?c), MODE(emotion in ?c's Moods) = ?e
    => dominantEmotion(?c, ?e)}\par
    \texttt{Character(?c), COUNT(DISTINCT emotion of ?c) = ?n
    => emotionalRange(?c, ?n)}\par
    \textit{(Cần aggregation; Cypher hoặc SPARQL + subquery argmax.)} \\
  \hline

  \texttt{B.7}
  & Costume Represents Mood
  & Costume xuất hiện cùng Mood $\geq 3$ lần thì thể hiện Mood đó.
  & \texttt{Costume(?co), Mood(?m),
    COUNT(?ra : hasAppearance(?ra,?co), hasAppearance(?ra,?m)) >= 3
    => represent(?co, ?m)}\par
    \textit{(Cần aggregation; SPARQL CONSTRUCT với
    GROUP BY ?co ?m HAVING COUNT(*) >= 3.)} \\
  \hline

  \texttt{C.1}
  & Expected Gender per Role Type
  & Quy ước nghệ thuật chèo: Đào và Mụ luôn là nữ; Kép và Lão luôn là nam.
  & \texttt{DaoCharacter rdfs:subClassOf
    [owl:onProperty charGender; owl:hasValue "Nữ"]}\par
    \texttt{KepCharacter rdfs:subClassOf
    [owl:onProperty charGender; owl:hasValue "Nam"]}\par
    \texttt{MuCharacter rdfs:subClassOf
    [owl:onProperty charGender; owl:hasValue "Nữ"]}\par
    \texttt{LaoCharacter rdfs:subClassOf
    [owl:onProperty charGender; owl:hasValue "Nam"]} \\
  \hline

  \texttt{C.2}
  & Tragic Heroine Archetype
  & Đào-Chín có $\geq 4$ cảm xúc gồm cảm xúc tiêu cực thì là nguyên mẫu nữ chính bi kịch.
  & \texttt{DaoCharacter(?c), subType(?c,"Chín"),
    COUNT(DISTINCT emotion of ?c) >= 4,
    EXISTS dark emotion of ?c IN \{Sad, Angry, Rage, Fear\}
    => archetype(?c, "Nữ chính bi kịch")}\par
    \textit{(Cần aggregation và filter exists.)} \\
  \hline

  \texttt{C.3}
  & Cross-Play Archetype Similarity
  & Hai nhân vật ở hai vở khác nhau cùng (mainType, subType) là cùng nguyên mẫu.
  & \texttt{Character(?c1), Character(?c2),
    roleType(?c1,?rt), roleType(?c2,?rt),
    subType(?c1,?st), subType(?c2,?st),
    hasCharacter(?p1,?c1), hasCharacter(?p2,?c2),
    notEqual(?c1,?c2), notEqual(?p1,?p2)
    => archetypeSimilar(?c1,?c2)} \\
  \hline

\end{longtable}
